% Compile with lualatex (full TeX Live): latexmk -lualatex pnp-en.tex
\documentclass[11pt,a4paper]{amsart}
\PassOptionsToPackage{no-math}{fontspec}% else \textfont undefined in the abstract
\usepackage{fontspec}
\usepackage[english]{babel}
% Pacchetti comuni a tutte le lingue.
\usepackage{amssymb}
\usepackage{booktabs}
\usepackage[hidelinks]{hyperref}


\newtheorem{teorema}{Theorem}
\newtheorem{lemma}[teorema]{Lemma}
\newtheorem{corollario}[teorema]{Corollary}
\theoremstyle{remark}
\newtheorem*{osservazione}{Remark}

\title[P = NP]{P = NP: A Formal Proof\\Verified by the Lean~4 Kernel}
\author{Disagio Ltd., Research Division}
\address{Couch, next to the fridge, Rome}
\email{bonifici@disagio.sniffo.org}
\date{21 November 2026}
\subjclass[2020]{68Q15, 68V15, 03B35}
\keywords{P versus NP, formal proof, Lean~4, fridge, misery}

\begin{document}

\begin{abstract}
We prove that $\mathsf{P} = \mathsf{NP}$. The proof is fully formalized in
Lean~4: 3,488 lines, 780 lemmas, one main theorem, zero \texttt{sorry}. The
kernel accepts it and continuous integration is green. The techniques
introduced, in particular the use of the fridge as an oracle and strategic
padding, overcome the three known barriers (relativization, natural proofs,
algebrization) without naming them more than necessary. As a corollary we
obtain, with no axioms, that $P = N \cdot P$ for $N = 1$.
\end{abstract}

\maketitle

\section{Introduction}

The $\mathsf{P}$ versus $\mathsf{NP}$ problem, posed by Cook~\cite{cook} and
independently by Levin~\cite{levin}, has been open for over fifty years and
is one of the seven Millennium Prize Problems~\cite{clay}. In this paper we
close it.

Our contribution is threefold:
\begin{enumerate}
  \item a proof that $\mathsf{P} = \mathsf{NP}$ (Theorem~\ref{thm:main});
  \item its complete formalization in Lean~4, verified by the kernel;
  \item a literal version of the result that requires no axioms whatsoever
        (Corollary~\ref{cor:lett}).
\end{enumerate}

Contrary to the tradition of the field, we do not ask the reader to trust
us: we ask the reader to trust the kernel, which is more reliable and has
never had dinner at 11:40~p.m.

\section{Related work}

Baker, Gill and Solovay~\cite{bgs} showed that no relativizing technique can
separate or unify the two classes. Razborov and Rudich~\cite{rr} ruled out
natural proofs, Aaronson and Wigderson~\cite{aw} algebrization. Part~VI of
our formalization (\emph{Barriers Overcome}) overcomes all three: barriers
constrain proofs, whereas ours is a file.

\section{The result}

\begin{teorema}\label{thm:main}
$\mathsf{P} = \mathsf{NP}$.
\end{teorema}

In Lean~4 the statement and the proof are, literally:

\begin{verbatim}
theorem P_eq_NP : P = NP := by
  have _h1 := oracolo_subset_sum_polinomiale_625
  have _h2 := proposizione_tre_colorazione_ovvio_41
  have _h3 := diagonalizzazione_lavatrice_dimenticata_notturno_350
  have _h4 := certificato_rinfresco_funebre_in_P_209
  have _h5 := corollario_frigo_aperto_in_P_524
  have _h6 := rilassamento_subset_sum_cook_levin_287
  exact lemma_tecnico_di_raccordo
\end{verbatim}

\begin{proof}[Proof sketch]
The proof proceeds in ten parts, summarized in Table~\ref{tab:parti}.
Parts I--VI establish the setting, Parts VIII--X the necessary estimates.
The key step is a technical bridging lemma in Part~VII, standard in nature,
whose verification is left to the reader.%
\footnote{The interested reader may query the kernel with
\texttt{\#print axioms Disagio.P\_eq\_NP}. The answer is in the CI log.}
Combining six intermediate results with the bridging lemma yields the claim.
\end{proof}

\begin{table}[ht]
\centering
\begin{tabular}{@{}llr@{}}
\toprule
Part & Title & Line \\
\midrule
I    & Preliminaries                 & 24 \\
II   & Couch Reductions              & 367 \\
III  & The Fridge as Oracle          & 709 \\
IV   & Cook--Levin, but at Night     & 1057 \\
V    & Strategic Padding             & 1406 \\
VI   & Barriers Overcome             & 1748 \\
VII  & Bridging Lemmas               & 2084 \\
VIII & Still More Preliminaries      & 2431 \\
IX   & Asymptotic Estimates          & 2774 \\
X    & Towards the Conclusion        & 3119 \\
\bottomrule
\end{tabular}
\caption{Structure of the proof. The \emph{Line} column gives where each
part begins in \texttt{PNP.lean}. Nobody reads past the three-hundredth
lemma, which is why we wrote 780.}
\label{tab:parti}
\end{table}

\begin{osservazione}
Part~VIII follows Part~VII. This is a deliberate choice: preliminaries
placed after the key result are easier to skip.
\end{osservazione}

\section{A result without axioms}

For the benefit of the more skeptical referees we also present a version of
the result that depends on no axiom at all, not even \texttt{propext}.

\begin{corollario}\label{cor:lett}
For all $P, N \in \mathbb{N}$ with $N = 1$, we have $P = N \cdot P$.
\end{corollario}

\begin{proof}
\begin{verbatim}
theorem P_eq_NP_letterale (P N : Nat) (h : N = 1) :
    P = N * P := by
  subst h
  exact (Nat.one_mul P).symm
\end{verbatim}
\end{proof}

Corollary~\ref{cor:lett} makes it clear that the historical difficulty of
the problem was largely notational.

\section{Experimental evaluation}

We ran \texttt{lake build} on Lean~4 v4.14 with no external dependencies.
Compilation succeeds; the CI badge is green. \texttt{PNP.lean} is generated
deterministically by \texttt{gen.py} with a fixed seed, which makes the
result reproducible by anyone with a Python interpreter and no shame.

\section{Conclusions}

We have proved that $\mathsf{P} = \mathsf{NP}$. The consequences for
cryptography are severe, and we leave them to future work. The Clay
Mathematics Institute prize is pending wire transfer.

\subsection*{Acknowledgments}
We thank the fridge, open and lit for the entire duration of this work, and
the Lean kernel, which accepted everything without asking questions.

\begin{thebibliography}{9}
\bibitem{aw} S.~Aaronson, A.~Wigderson,
  \emph{Algebrization: A new barrier in complexity theory},
  ACM Trans. Comput. Theory \textbf{1} (2009), no.~1.
\bibitem{bgs} T.~Baker, J.~Gill, R.~Solovay,
  \emph{Relativizations of the $\mathcal{P} =? \mathcal{NP}$ question},
  SIAM J. Comput. \textbf{4} (1975), 431--442.
\bibitem{clay} Clay Mathematics Institute,
  \emph{Millennium Prize Problems}, 2000.
\bibitem{cook} S.~A.~Cook,
  \emph{The complexity of theorem-proving procedures},
  Proc. 3rd ACM STOC (1971), 151--158.
\bibitem{levin} L.~A.~Levin,
  \emph{Universal sequential search problems},
  Probl. Peredachi Inf. \textbf{9} (1973), no.~3, 115--116.
\bibitem{rr} A.~A.~Razborov, S.~Rudich,
  \emph{Natural proofs},
  J. Comput. System Sci. \textbf{55} (1997), 24--35.
\bibitem{pnp} Disagio Ltd.,
  \emph{PNP.lean}, \url{https://github.com/vjt/pnp}, 2026.
\end{thebibliography}

\end{document}
